% GENERATED FILE. Edit canonical lessons, then run npm run generate:books.
% canonical-source-hash: fnv1a64:2fbdb66498ab9a77
% canonical-lessons: TE-C438-prayanikudu, TE-C438-prayanincu, TE-C438-gamyam, TE-C438-margam, TE-C438-adugupettu

\chapter{Setting Out}
\label{ch:te-setting-out}
\hlchaptermodality{438}

\begin{chapteropening}
\textbf{By the end of this chapter:} \emph{I can say a passenger, to travel, a destination, a route and to set foot in a place.}

Complete the last lesson of chapter 438: I can say a passenger, to travel, a destination, a route and to set foot in a place.
\end{chapteropening}

\section[\texorpdfstring{\te{ప్రయాణికుడు} (prayāṇikuḍu)}{ప్రయాణికుడు (prayāṇikuḍu)}]{\te{ప్రయాణికుడు} (prayāṇikuḍu) — a passenger; a traveller}

\label{lesson:TE-C438-prayanikudu}



\begin{quote}
Before the new one: say the Telugu for a watchman, then the Telugu for a milkman.
\end{quote}

\subsection*{You'll want to know: \te{ప్రయాణికుడు}}
\textbf{\te{ప్రయాణికుడు}} — \emph{prayāṇikuḍu} — \textquotedblleft{}a passenger; a traveller\textquotedblright{}.

\begin{quote}
\te{రైలులో} \te{చాలా} \te{మంది} \te{ప్రయాణికులు} \te{ఉన్నారు}.
\end{quote}

\begin{quote}
\emph{railulō cālā mandi prayāṇikulu unnāru.}
\end{quote}

\begin{quote}
There were a lot of passengers on the train.
\end{quote}

Plural: \te{ప్రయాణికులు}, \emph{prayāṇikulu}.

The ending -\te{ఇకుడు} makes a person out of a noun: \te{ప్రయాణం}, a journey, gives \te{ప్రయాణికుడు}, one who is on a journey. A woman passenger is \te{ప్రయాణికురాలు}.

\begin{cousinweb}
\textbf{\te{ప్రయాణికుడు}} (\emph{prayāṇikuḍu}) is built on \te{ప్రయాణం}, \textquotedblleft{}a journey\textquotedblright{}, which is borrowed from Sanskrit \emph{prayāṇa}, \textquotedblleft{}setting out, a journey\textquotedblright{} (CDIAL 8732). The person ending -\te{ఇకుడు}, \textquotedblleft{}one who …\textquotedblright{}, makes the traveller.
\end{cousinweb}

\begin{grammarlens}[title={What you've built}]
One more word for getting around.
\end{grammarlens}

\subsection*{Your turn}
\begin{itemize}
\raggedright
  \item \emph{Say it:} \emph{prayāṇikuḍu}
  \item \emph{Say it:} \emph{prayāṇikuḍu}, once more
  \item \emph{Say it:} \emph{railulō cālā mandi prayāṇikulu unnāru.} — There were a lot of passengers on the train.
  \item \emph{Recall:} say the Telugu for a watchman, then the Telugu for a milkman, then say \emph{prayāṇikuḍu} again
\end{itemize}

\begin{tcolorbox}[breakable,colback=teal!4,colframe=teal!35!black,title={Before you move on}]
What does \te{ప్రయాణికుడు} mean? (\textquotedblleft{}A passenger; a traveller\textquotedblright{}.) Say it once more.
\end{tcolorbox}
\section[\texorpdfstring{\te{ప్రయాణించు} (prayāṇiñcu)}{ప్రయాణించు (prayāṇiñcu)}]{\te{ప్రయాణించు} (prayāṇiñcu) — to travel}

\label{lesson:TE-C438-prayanincu}



\begin{quote}
Before the new one: say the Telugu for a milkman, then the Telugu for a passenger.
\end{quote}

\subsection*{You'll want to know: \te{ప్రయాణించు}}
\textbf{\te{ప్రయాణించు}} — \emph{prayāṇiñcu} — \textquotedblleft{}to travel\textquotedblright{}.

\begin{quote}
\te{మేము} \te{రాత్రంతా} \te{బస్సులో} \te{ప్రయాణించాం}.
\end{quote}

\begin{quote}
\emph{mēmu rātrantā bassulō prayāṇiñcāṁ.}
\end{quote}

\begin{quote}
We travelled by bus all night.
\end{quote}

Past (I …): \te{ప్రయాణించాను}, \emph{prayāṇiñcānu}.

The ending -\te{ఇంచు} turns a Sanskrit noun into a Telugu verb: \te{ప్రయాణం}, a journey, gives \te{ప్రయాణించు}, to make a journey. The means of travel takes -\te{లో}: \te{రైలులో}, \te{బస్సులో}.

\begin{cousinweb}
\textbf{\te{ప్రయాణించు}} (\emph{prayāṇiñcu}) is built from \te{ప్రయాణం}, borrowed from Sanskrit \emph{prayāṇa}, \textquotedblleft{}setting out, a journey\textquotedblright{} (CDIAL 8732), with the Telugu verb-making ending -\te{ఇంచు}.
\end{cousinweb}

\begin{grammarlens}[title={What you've built}]
One more word for getting around.
\end{grammarlens}

\subsection*{Your turn}
\begin{itemize}
\raggedright
  \item \emph{Say it:} \emph{prayāṇiñcu}
  \item \emph{Say it:} \emph{prayāṇiñcu}, once more
  \item \emph{Say it:} \emph{mēmu rātrantā bassulō prayāṇiñcāṁ.} — We travelled by bus all night.
  \item \emph{Recall:} say the Telugu for a milkman, then the Telugu for a passenger, then say \emph{prayāṇiñcu} again
\end{itemize}

\begin{tcolorbox}[breakable,colback=teal!4,colframe=teal!35!black,title={Before you move on}]
What does \te{ప్రయాణించు} mean? (\textquotedblleft{}To travel\textquotedblright{}.) Say it once more.
\end{tcolorbox}
\section[\texorpdfstring{\te{గమ్యం} (gamyaṁ)}{గమ్యం (gamyaṁ)}]{\te{గమ్యం} (gamyaṁ) — a destination}

\label{lesson:TE-C438-gamyam}



\begin{quote}
Before the new one: say the Telugu for a passenger, then the Telugu for to travel.
\end{quote}

\subsection*{You'll want to know: \te{గమ్యం}}
\textbf{\te{గమ్యం}} — \emph{gamyaṁ} — \textquotedblleft{}a destination\textquotedblright{}.

\begin{quote}
\te{మా} \te{గమ్యం} \te{ఇంకా} \te{చాలా} \te{దూరంలో} \te{ఉంది}.
\end{quote}

\begin{quote}
\emph{mā gamyaṁ iṅkā cālā dūraṁlō undi.}
\end{quote}

\begin{quote}
Our destination is still a long way off.
\end{quote}

A formal word; in everyday talk people ask where you are going: \te{ఎక్కడికి} \te{వెళ్తున్నారు}?

\begin{cousinweb}
\textbf{\te{గమ్యం}} (\emph{gamyaṁ}) is borrowed from Sanskrit \emph{gamya}, \textquotedblleft{}to be gone to, reachable\textquotedblright{}, from the root \emph{gam}, \textquotedblleft{}to go\textquotedblright{} (CDIAL 4032). The place one can go to became the place one is going to.
\end{cousinweb}

\begin{grammarlens}[title={What you've built}]
One more word for getting around.
\end{grammarlens}

\subsection*{Your turn}
\begin{itemize}
\raggedright
  \item \emph{Say it:} \emph{gamyaṁ}
  \item \emph{Say it:} \emph{gamyaṁ}, once more
  \item \emph{Say it:} \emph{mā gamyaṁ iṅkā cālā dūraṁlō undi.} — Our destination is still a long way off.
  \item \emph{Recall:} say the Telugu for a passenger, then the Telugu for to travel, then say \emph{gamyaṁ} again
\end{itemize}

\begin{tcolorbox}[breakable,colback=teal!4,colframe=teal!35!black,title={Before you move on}]
What does \te{గమ్యం} mean? (\textquotedblleft{}A destination\textquotedblright{}.) Say it once more.
\end{tcolorbox}
\section[\texorpdfstring{\te{మార్గం} (mārgaṁ)}{మార్గం (mārgaṁ)}]{\te{మార్గం} (mārgaṁ) — a route; a way}

\label{lesson:TE-C438-margam}



\begin{quote}
Before the new one: say the Telugu for to travel, then the Telugu for a destination.
\end{quote}

\subsection*{You'll want to know: \te{మార్గం}}
\textbf{\te{మార్గం}} — \emph{mārgaṁ} — \textquotedblleft{}a route; a way\textquotedblright{}.

\begin{quote}
\te{ఆ} \te{మార్గంలో} \te{బస్సులు} \te{తక్కువగా} \te{ఉంటాయి}.
\end{quote}

\begin{quote}
\emph{ā mārgaṁlō bassulu takkuvagā uṇṭāyi.}
\end{quote}

\begin{quote}
There are few buses on that route.
\end{quote}

\te{దారి} is the everyday word for a way or a road; \te{మార్గం} is a little more formal, and it also means a way of doing something.

\begin{cousinweb}
\textbf{\te{మార్గం}} (\emph{mārgaṁ}) is borrowed from Sanskrit \emph{mārga}, \textquotedblleft{}a track, a road\textquotedblright{} (CDIAL 10071). CDIAL notes that it first meant the trail of a wild animal.
\end{cousinweb}

\begin{grammarlens}[title={What you've built}]
One more word for getting around.
\end{grammarlens}

\subsection*{Your turn}
\begin{itemize}
\raggedright
  \item \emph{Say it:} \emph{mārgaṁ}
  \item \emph{Say it:} \emph{mārgaṁ}, once more
  \item \emph{Say it:} \emph{ā mārgaṁlō bassulu takkuvagā uṇṭāyi.} — There are few buses on that route.
  \item \emph{Recall:} say the Telugu for to travel, then the Telugu for a destination, then say \emph{mārgaṁ} again
\end{itemize}

\begin{tcolorbox}[breakable,colback=teal!4,colframe=teal!35!black,title={Before you move on}]
What does \te{మార్గం} mean? (\textquotedblleft{}A route; a way\textquotedblright{}.) Say it once more.
\end{tcolorbox}
\section[\texorpdfstring{\te{అడుగుపెట్టు} (aḍugupeṭṭu)}{అడుగుపెట్టు (aḍugupeṭṭu)}]{\te{అడుగుపెట్టు} (aḍugupeṭṭu) — to set foot in; to arrive}

\label{lesson:TE-C438-adugupettu}



\begin{quote}
Before the new one: say the Telugu for a destination, then the Telugu for a route.
\end{quote}

\subsection*{You'll want to know: \te{అడుగుపెట్టు}}
\textbf{\te{అడుగుపెట్టు}} — \emph{aḍugupeṭṭu} — \textquotedblleft{}to set foot in; to arrive\textquotedblright{}.

\begin{quote}
\te{మేము} \te{రాత్రి} \te{పది} \te{గంటలకు} \te{కొత్త} \te{నగరంలో} \te{అడుగుపెట్టాం}.
\end{quote}

\begin{quote}
\emph{mēmu rātri padi gaṇṭalaku kotta nagaraṁlō aḍugupeṭṭāṁ.}
\end{quote}

\begin{quote}
We set foot in the new city at ten at night.
\end{quote}

Past (I …): \te{అడుగుపెట్టాను}, \emph{aḍugupeṭṭānu}.

\te{అడుగు}, a step, + \te{పెట్టు}, to put: to put one's foot into a place, to arrive. The place takes -\te{లో}: \te{నగరంలో} \te{అడుగుపెట్టాను}. \te{అడుగు}, to ask, is another word.

\begin{cousinweb}
\textbf{\te{అడుగుపెట్టు}} (\emph{aḍugupeṭṭu}) is native Dravidian: \te{అడుగు}, \textquotedblleft{}a foot, a footstep, a step\textquotedblright{} (DEDR 72, one word with Tamil \emph{aṭi}, \textquotedblleft{}foot\textquotedblright{}), and \te{పెట్టు}, \textquotedblleft{}to put, to place\textquotedblright{} (DEDR 4390). \te{అడుగు}, \textquotedblleft{}to ask\textquotedblright{}, is a separate word (DEDR 81).
\end{cousinweb}

\begin{grammarlens}[title={What you've built}]
One more word for getting around.
\end{grammarlens}

\subsection*{Your turn}
\begin{itemize}
\raggedright
  \item \emph{Say it:} \emph{aḍugupeṭṭu}
  \item \emph{Say it:} \emph{aḍugupeṭṭu}, once more
  \item \emph{Say it:} \emph{mēmu rātri padi gaṇṭalaku kotta nagaraṁlō aḍugupeṭṭāṁ.} — We set foot in the new city at ten at night.
  \item \emph{Recall:} say the Telugu for a destination, then the Telugu for a route, then say \emph{aḍugupeṭṭu} again
\end{itemize}

\begin{tcolorbox}[breakable,colback=teal!4,colframe=teal!35!black,title={Before you move on}]
What does \te{అడుగుపెట్టు} mean? (\textquotedblleft{}To set foot in; to arrive\textquotedblright{}.) Say it once more.
\end{tcolorbox}
